\section{SUMMARY AND CONCLUSION}

\subsection{Summary}

This study developed and evaluated \textit{Pinoy Speak}, a
deployed full-stack system for tracking, organizing, and
explaining informal Filipino and Taglish slang. The system
addresses the limitations of static dictionaries by
continuously collecting public online text, identifying
candidate slang terms, enriching verified entries, and
presenting them through user-facing learning tools.

The current implementation archives 175,481 records from
Reddit, YouTube comments, and user chat interactions. It
maintains approximately 104 verified slang entries composed
of seed entries and dynamically discovered terms. The system
uses a multi-stage NLP pipeline that combines calamanCy
tokenization, dictionary gating, subword novelty signals,
FastText-based similarity support, temporal burstiness
analysis, morphological filtering, and context-window
disambiguation. These components allow the system to detect
both novel out-of-vocabulary slang and known words that may
be undergoing semantic shift.

\subsection{Conclusion}

The evaluation shows that the system is promising but still
requires iterative human validation. The 79-entry lexicon
validation produced strong ratings for definition correctness
(4.83/5) and example naturalness (4.75/5). Most item-level
judgments were marked correct or acceptable, but validators
also identified entries that need revision, especially cases
where the word may be ordinary Filipino vocabulary, borrowed
English, or a context-dependent expression rather than clear
slang. This confirms that human review remains necessary for
maintaining a reliable living dictionary.

The 36-response consented beta usability evaluation showed
positive user acceptance. The Dictionary received the highest
module usability rating at 9.36 out of 10, followed by the
Tutor at 9.33, Navigation Flow at 9.14, Analyzer or Translator
at 9.00, Top Slang at 8.85, and Dashboard or Word Popularity
at 8.58. In addition, 86.11\% of respondents reported
learning at least one new slang word or expression from the
system.

Overall, the study shows that informal Filipino and Taglish
slang can be modeled as dynamic linguistic data rather than
treated merely as noise. By integrating corpus-based trend
analysis, contextual disambiguation, LLM-assisted enrichment,
human lexicon validation, and an educational web interface,
\textit{Pinoy Speak} provides an application-oriented
contribution to Filipino NLP and digital language learning.

\subsection{Recommendations}

Although the implemented system achieved its main objectives,
several improvements are recommended for future work.

First, the corpus should be expanded to include more regional
Filipino languages, dialectal varieties, and platform sources.
The current corpus is useful for detecting slang from selected
high-activity online communities, but it does not represent
all Filipino speakers or all regional forms of informal
language.

Second, lexicon validation should be expanded. The current
validation provided three ratings per term, which is useful
for initial checking, but future versions should include more
annotators and clearer annotation guidelines. Priority should
be given to entries flagged by validators, including terms
with unclear slang status, incorrect definitions, weak
examples, or questionable formation labels.

Third, the system should improve its handling of ambiguous
and borrowed words. Validator comments showed that terms such
as \textit{basic}, \textit{dead}, \textit{legit},
\textit{savage}, \textit{shook}, and \textit{shookt} require
stronger explanation of semantic shift, borrowing, and local
usage.

Fourth, the user interface should refine chart interactions,
dictionary navigation, and tutor activities. Usability
feedback suggested clearer graph controls, better loading
indicators, more natural tutor responses, and separate quiz
or flashcard features.

Finally, future evaluation should include more formal
computational metrics such as precision, recall,
false-positive rate, false-negative rate, and
inter-annotator agreement. These measures would strengthen
the evidence for system accuracy beyond usability ratings and
initial human validation.

\subsection{Final Remarks}

\textit{Pinoy Speak} demonstrates that Filipino internet slang
can be documented through a combination of computational
detection, human validation, and interactive presentation.
The system is not intended to replace expert linguistic
judgment. Rather, it serves as a practical tool for tracking
emerging expressions, helping users understand context, and
supporting future work in low-resource Filipino NLP.
